% ============================================================
%  Chapter 02: Sociology — The Discipline: Emergence, 
%  Theoretical Foundations, Common Sense & Scope
% ============================================================

\chapter{Emergence, Theoretical Foundations, Common Sense \& Scope of Sociology}

\section{Emergence of Sociology, Positivism \& Structural Functionalism}

\subsection{Approach to Sociology Answer Writing (UPSC Mains)}

\begin{DataFact}
\begin{itemize}
    \item \textbf{UPSC Civil Services Mains, Sociology Optional:} Total time 3 hours; 28 questions given, 19 to be attempted .
    \item Of the 19 attempted: 13 are 10-mark questions and 6 are 20-mark questions .
    \item \textbf{Time available:} Approximately 8 minutes for a 10-mark question and 11--12 minutes for a 20-mark question---an average of about 10 minutes per question across the paper .
    \item In the exam hall, aspirants compete against time, not against other candidates---comparing preparation with peers (e.g., on group chats) creates unnecessary anxiety, which should be controlled .
\end{itemize}
\end{DataFact}

\paragraph{The Four-Step Method for Any Question}
Whether a question carries 10 or 20 marks, the approach (not the content) comes first . The same four steps apply to every question :
\begin{enumerate}
    \item \textbf{Step 1 --- Identify the keywords} in the question, before reading the question as a whole .
    \item \textbf{Step 2 --- Define each identified keyword} precisely .
    \item \textbf{Step 3 --- Read the full question} only after defining the keywords, to see the relationship/pattern being asked about .
    \item \textbf{Step 4 --- Identify the pattern:} Most questions follow a causation pattern (``How did A lead to B'') . Causation means a cause--effect relationship: the cause is the independent variable and the effect is the dependent variable (e.g., $\text{smoking} = \text{independent variable}$, $\text{cancer} = \text{dependent variable}$) .
    \item Where a question has no explicit causation (e.g., ``Define A'' or ``What factors lead to A''), the same logic applies: A becomes the dependent variable and its factors become the independent variables .
\end{enumerate}

\begin{PYQLink}
\textbf{Worked example (asked by UPSC in a previous year):} \textit{``How did the intellectual forces lead to the emergence of sociology? Discuss.''} (10 marks) 
\begin{itemize}
    \item \textbf{Keywords identified:} `Intellectual forces' and `emergence of sociology' .
    \item \textbf{Keyword 1 defined:} Intellectual forces are forces characterised by a change in ideas or thinking .
    \item \textbf{Pattern:} Causation---intellectual forces (independent variable) led to the emergence of sociology (dependent variable) .
\end{itemize}
\end{PYQLink}

\paragraph{Answer Length \& Diagram Use}
\begin{itemize}
    \item A 10-mark answer should run to roughly 150 words---about one page (both sides), i.e., two written sides .
    \item A short, clean diagram helps orient both the writer and the examiner :
    $$\text{Middle Ages} \longrightarrow \text{Intellectual Forces} \longrightarrow \text{Modernity} \longrightarrow \text{Conservative/Radical Reaction} \longrightarrow \text{Emergence of Sociology}$$
    UPSC does not award separate marks for diagrams, but they demonstrate clarity of thought and individuality .
\end{itemize}

\begin{CriticalPoint}
\begin{itemize}
    \item Write in keywords, not long explanatory paragraphs---keywords are the aspirant's real weapon in a 150-word answer .
    \item Examiners evaluate a very large number of scripts in limited time and are themselves specialists evaluating aspirants, not sociology postgraduates---long paragraphs risk irritating the examiner and losing marks rather than gaining them .
\end{itemize}
\end{CriticalPoint}

\begin{QuickRevision}
\begin{itemize}
    \item \textbf{UPSC Sociology Optional:} 3 hrs, 19/28 questions, $13 \times 10\text{-mark} + 6 \times 20\text{-mark}$, $\sim 10\text{ min/question}$ average .
    \item \textbf{Four-step method:} Identify keywords $\rightarrow$ define keywords $\rightarrow$ read the question $\rightarrow$ spot the pattern (usually causation: independent variable = cause, dependent variable = effect) .
    \item 10-mark answer $\approx 150$ words / 2 sides; write in precise keywords, not long prose; diagrams aid clarity though not separately marked .
\end{itemize}
\end{QuickRevision}

\subsection{Intellectual Forces Leading to the Emergence of Sociology}

\begin{KeyConcept}
\begin{itemize}
    \item \textbf{Intellectual forces:} Forces characterised by a change in ideas or thinking (as distinct from material/economic forces) .
    \item \textbf{Renaissance:} Its biggest contribution was the ``invention of man and his potential''---as the authority of God/religion declined or was questioned, man as a rational, capable individual was `discovered' .
    \begin{itemize}
        \item Key terms for writing: \textbf{humanism, individualism, secularism} .
    \end{itemize}
    \item \textbf{Scientific Renaissance:} The manifestation of man's discovered potential in exploring the natural world, through observation and experimentation (e.g., Galileo Galilei) .
    \item \textbf{Enlightenment:} The same spirit of curiosity and reason, now turned toward the social world; Enlightenment scholars, influenced by the scientific revolution, sought a scientific/rational form of governance for society .
    \item \textbf{French Revolution:} Treated as an effect/outcome of Enlightenment ideas .
    \item \textbf{Counter-Enlightenment:} The conservative reaction that resisted the Enlightenment/French Revolution's changes; resistance to an idea is itself a form of intellectual change, so Counter-Enlightenment also counts as an intellectual force .
\end{itemize}
\end{KeyConcept}

\begin{CriticalPoint}
\begin{itemize}
    \item \textbf{Industrial Revolution should NOT be listed as an intellectual force}---it is a material force, not a change in ideas/thinking .
    \item \textbf{Analogy:} Education is an idea (intellectual force); a textbook is a concrete/material object---part of education but not itself an intellectual force .
    \item Conservative reaction to the French Revolution and Enlightenment gave rise to \textbf{`French sociology'}, which was conservative in character---it favoured reform, not revolution .
\end{itemize}
\end{CriticalPoint}

\begin{QuickRevision}
\begin{itemize}
    \item \textbf{Sequence:} Renaissance $\rightarrow$ Scientific Renaissance $\rightarrow$ Enlightenment $\rightarrow$ French Revolution $\rightarrow$ Counter-Enlightenment (conservative reaction) $\rightarrow$ Emergence of Sociology .
    \item \textbf{Renaissance keywords:} Humanism, individualism, secularism .
    \item Industrial Revolution = material force, \textbf{NOT} an intellectual force---common error to avoid in writing .
\end{itemize}
\end{QuickRevision}

\subsection{Auguste Comte, `Intellectual Disorder' and Positivism}
Auguste Comte was a positivist scholar; the duty of a positivist scholar is to uncover the laws that govern the social world, just as natural laws (e.g., day following night) govern the natural world .

\begin{KeyConcept}
\textbf{Intellectual disorder (Comte):} Even though society had entered the scientific age (around 1800), people continued to explain the world in religious and metaphysical (abstract) terms rather than scientific terms---this persistence of pre-scientific thought amid a scientific age is what Comte called intellectual disorder .
\begin{itemize}
    \item Comte, writing in the context of the French Revolution and Enlightenment, observed visible disorder around him (executions, revolts) and, as a self-styled `scientist of society', claimed he could identify why disorder existed and predict when order would be restored .
    \item \textbf{Reason for disorder, per Comte:} Religious and metaphysical explanations still competed with scientific explanations in the popular mind---order would return only once one single form of knowledge (scientific knowledge) was accepted by everyone without dispute .
\end{itemize}
\end{KeyConcept}

\begin{CriticalPoint}
\begin{itemize}
    \item Comte did not `undermine' or condemn religion as such---he was himself a conservative who did not like religion or monarchy being challenged; he was scientifically describing the co-existence of religious/metaphysical and scientific thought as the cause of disorder, not attacking religion .
    \item \textbf{Modern illustration cited:} Ashis Nandy and T.~N. Madan have argued that although India is constitutionally secular, Indians remain deeply religious in practice---offered as an example of persisting non-scientific thought .
\end{itemize}
\end{CriticalPoint}

\paragraph{Law of (Three) Stages}
Comte held that every society passes through three stages, analogous to the stages of a human life (infancy, adulthood, old age) :
\begin{enumerate}
    \item \textbf{Theological stage (up to c.~1300 AD):} Everything is explained in religious terms (e.g., `obey the priestly class to reach heaven'); corresponds to the medieval/dark ages, before the Renaissance .
    \item \textbf{Metaphysical stage (c.~1300--1800 AD):} An abstract, in-between stage: society has moved past pure religious explanation but has not yet reached concrete scientific explanation . The intellectual forces from Renaissance to Counter-Enlightenment unfolded within this $\sim 500$-year stage .
    \item \textbf{Scientific stage (from c.~1800 onward):} Everything is explained through scientific facts rather than religion or abstraction .
\end{enumerate}

\begin{ExampleBox}
\begin{itemize}
    \item \textbf{Rain example across the three stages:}
    \begin{itemize}
        \item \textit{Theological stage:} Perform rituals/havan to bring rain .
        \item \textit{Metaphysical stage:} Reject ritual but simply wait, without yet knowing the scientific cause .
        \item \textit{Scientific stage:} Spraying silver iodide ($\text{AgI}$) on clouds to induce artificial rain .
    \end{itemize}
    \item \textbf{`Man the conqueror'} is the idea of the scientific stage: the same man once thought `cursed' by asteroids sent by God can now, in the scientific stage, deflect an asteroid .
\end{itemize}
\end{ExampleBox}

\paragraph{Comte's Vision for Sociology \& Positivism}
\begin{itemize}
    \item Comte believed sociology could: (1) help societies advance into the scientific stage; (2) help society understand the reasons for disorder; (3) help study the laws governing social order and social disorder .
    \item Positivism, per Comte, is the philosophy that can help discover these laws; its underlying assumption is that society can be studied scientifically .
\end{itemize}

\begin{KeyConcept}
\begin{itemize}
    \item \textbf{Comte's thesis:} ``Much as the natural world is governed by laws of physics, so does society [have its own laws]'' .
    \item The biggest challenge for a sociologist, unlike a doctor (who can see/dissect the human body) or a geologist (who can see/cut rocks), is that society itself cannot be directly seen---only fragments of it (individuals, organisations) are visible .
\end{itemize}
\end{KeyConcept}

\begin{CriticalPoint}
\begin{itemize}
    \item Because sociology studies something not directly visible, Comte called sociology the \textbf{`Queen of Sciences'}---uncovering laws of an invisible social world is a tougher challenge than studying a visible natural object .
    \item Critics have called positivism a new form of religion, and positivists a new form of clergy---flagged as a critical debate .
    \item \textbf{Predictive science:} Besides uncovering laws, positivists claim to predict the future of the social world, just as natural scientists predict natural events .
\end{itemize}
\end{CriticalPoint}

\begin{GovtPolicy}
\begin{itemize}
    \item \textbf{Marxism as a positivist philosophy:} It too claims to uncover laws of the social world and to predict---Marx predicted that revolution is inevitable, that the working class will overthrow the bourgeoisie and establish socialism, followed by communism .
    \item \textbf{Personal example (Teacher's experience at DRDO, Jodhpur):} Working on a repetitive, highly specialised alloy-development task for months raised the question of whether this was Marx's `alienation'---doing a specialised, repetitive task can make one feel like ``a cog in the machine'' .
    \item \textbf{Contrast offered:} The IAS is considered the one service (besides the IFS) that avoids alienation, because officers rotate across very different departments (e.g., Education one year, Mines the next) .
\end{itemize}
\end{GovtPolicy}

\begin{QuickRevision}
\begin{itemize}
    \item \textbf{Intellectual disorder:} Persistence of religious/metaphysical thought despite society having entered the scientific age . Comte diagnosed this co-existence as the cause of disorder .
    \item \textbf{Law of Three Stages:} Theological ($-\text{1300}$) $\rightarrow$ Metaphysical ($1300\text{--}1800$) $\rightarrow$ Scientific ($1800-$) .
    \item \textbf{Comte:} Sociology should study laws of social order/disorder; positivism assumes society can be studied scientifically and is predictive .
    \item \textbf{`Queen of Sciences':} Sociology is so named because society cannot be seen as a whole .
    \item \textbf{Marxism:} Positivist philosophy predicting inevitable revolution $\rightarrow$ socialism $\rightarrow$ communism .
\end{itemize}
\end{QuickRevision}

\subsection{Precursors to Sociology}

\begin{KeyConcept}
A \textbf{precursor} is a discipline that existed prior to sociology and addressed, in a partial way, questions that sociology would later take up as a full discipline .
\begin{enumerate}
    \item \textbf{Philosophy of History:} Held that society moves in a linear fashion, from a simple to a complex stage .
    \item \textbf{Biological Theory of Evolution:} Compared society to a human body: just as the body grows from infancy to adulthood through multiple organs (heart, liver, kidney, brain---all equally important), society too consists of multiple, equally important parts (economy, religion, family, media, state, law) .
    \begin{itemize}
        \item \textbf{Herbert Spencer} is credited as the scholar who first compared human body and society in this way (later termed the \textbf{`organic analogy'}) .
    \end{itemize}
    \item \textbf{Survey of Social Conditions:} Through survey-based study, it was discovered for the first time that poverty is not a creation of God/fate but is socially constructed---specifically, a structural condition arising from industrialisation (e.g., slum development) .
\end{enumerate}
\end{KeyConcept}

\begin{ExampleBox}
\begin{itemize}
    \item \textbf{Illustration:} Dalits are, on data, relatively poorer than upper-caste groups---not because of any `natural' deficiency, but because of caste as a social factor (itself rooted in religion), showing poverty to be structurally caused rather than natural .
    \item \textbf{Hymn reflecting older pre-survey view:} \textit{``The rich man in his castle, the poor man at his gate, God made them high and lowly, and ordered their estate''} .
    \item \textbf{Modern policy classification:} Government poverty classifications such as \textit{Antyodaya Anna Yojana} and BPL/APL categories---poverty is now studied and classified, not attributed to divine will .
    \item \textbf{Exam note:} In an answer on the emergence of sociology, combining the intellectual forces with these three precursors accounts for roughly 60--70 words of the 150 words expected .
\end{itemize}
\end{ExampleBox}

\begin{QuickRevision}
\begin{itemize}
    \item \textbf{Three precursors:} Philosophy of History (simple $\rightarrow$ complex), Biological Theory of Evolution (organic analogy---Herbert Spencer), Survey of Social Conditions (poverty as structurally/socially constructed) .
    \item Combine intellectual forces + 3 precursors for a complete 10-marker on the emergence of sociology .
\end{itemize}
\end{QuickRevision}

\subsection{Doubts Addressed (Day 1)}

\begin{DataFact}
\begin{itemize}
    \item \textbf{American sociology} (covered under Unit 4) includes three syllabus scholars: Talcott Parsons, Robert K. Merton, and George Herbert Mead .
    \item \textbf{Nobel Prize in Sociology:} Only one sociologist has ever won a Nobel Prize---\textbf{Jane Addams}, a female sociologist who worked on poverty and slums in America (awarded the Nobel Peace Prize) .
\end{itemize}
\end{DataFact}

\begin{itemize}
    \item \textbf{Philosophy of History (re-explained):} History/society moves from simple to complex---e.g., France once had one dominant religion and one form of family; today, with religious/social diversity, the state must be secular (equal distance from all religions) to survive .
    \item \textbf{Social thought vs. sociology:} Social thought is a very general term (how anyone thinks about society), not a discipline in itself; sociology is a formal, scientific discipline .
    \item \textbf{Why Comte called disorder `intellectual':} If, even after enlightenment/scientific progress, society still relies on religious or caste-based belief for legitimacy, that persistence of pre-scientific thinking is itself the disorder---just as constitutional morality requires everyone to actually believe in the constitution, not merely have it exist on paper .
    \item \textbf{Why early sociologists wanted sociology to be scientific:} The dominant mode of thought was scientific observation and experimentation (Galileo, Kepler, Copernicus) and the rational, doubting man; sociology's founders wanted the same rigour to uncover the laws governing society .
\end{itemize}

\begin{CriticalPoint}
\begin{itemize}
    \item \textbf{Metanarrative:} Grand, totalising belief systems (``\textit{badi badi baatein}'')---e.g., Marxism's promise of moving the world toward egalitarianism . Such ideologies are increasingly seen as having `failed' in practice: even self-declared communist parties today often support positions that do not strictly follow their own ideology, showing ground reality has diverged from the grand ideological narrative .
    \item \textbf{Nuclear family as a concept:} \textbf{A.~M. Shah} (sociologist, Delhi School of Economics, colleague of André Béteille) argued that the concept of the `nuclear family', in its strict Western sense of complete disconnection from extended kin, does not really exist in India---even when Indian children live away from parents, daily contact with parents typically continues .
\end{itemize}
\end{CriticalPoint}

\begin{QuickRevision}
\begin{itemize}
    \item American sociology scholars: Parsons, Merton, Mead .
    \item Only Nobel laureate sociologist: Jane Addams (poverty/slum work, USA) .
    \item Metanarrative = grand totalising ideological belief that has lost practical traction .
    \item A.~M. Shah: Western-style disconnected `nuclear family' does not truly exist in India .
\end{itemize}
\end{QuickRevision}

\subsection{Structural Functionalism (Introduction)}

\begin{CriticalPoint}
\begin{itemize}
    \item \textbf{Synonymous terms:} Structural functionalism $\equiv$ functionalism $\equiv$ structural theory $\equiv$ consensus theory $\equiv$ consensus structuralism---all refer to the same overarching theory .
    \item \textbf{Key scholars:} Herbert Spencer, Auguste Comte, Émile Durkheim .
    \item \textbf{Core idea:} Society is analogous to the human body . The body has multiple organs (brain, heart, lungs, liver, kidney), each performing a function needed for survival; society similarly has multiple parts/organs (religion, economy, law, family, state, media), each performing a vital function for social survival .
\end{itemize}
\end{CriticalPoint}

\begin{ExampleBox}
\begin{itemize}
    \item \textbf{Function of religion:} Provides norms/ethics for a moral life . Example---the Ten Commandments in Christianity; the five precepts (\textit{Panchsheel}) in Buddhism (no killing, no stealing, no lying, no intoxicants, no sexual misconduct); Jain teachings similarly .
    \item \textbf{Classroom illustration:} If all students share the same religious norms (e.g., Buddhist norms against alcohol), order is maintained; introducing one member with very different norms creates conflict---shared norms sustain social order .
    \item \textbf{Social equilibrium:} The state of maintained order that results when all social organs function properly---borrowed from the biological concept of \textbf{homeostasis} .
    \item \textbf{Function of law:} Prevents deviance and activates when norms/order are broken . \textit{Personal example:} Waiting at a red light while another rider jumps it; the rule-breaker is later caught by police, restoring faith that law maintains order .
    \item \textbf{Pathology:} When an organ of society malfunctions or stops functioning .
    \item \textbf{Cure for pathology:} Social reforms, not revolution---revolution is compared to an organ transplant (a last resort, not the first step) .
\end{itemize}
\end{ExampleBox}

\begin{CriticalPoint}
\begin{itemize}
    \item If organs fail altogether, society becomes `disabled'---disorder results .
    \item \textbf{Economic malfunctions:} Recession, inflation, hyperinflation, stagflation .
    \item \textbf{Lockdown illustration:} The COVID-19 lockdown is framed as the state/law `organ' functioning to prevent greater systemic collapse (more deaths) .
    \item \textbf{French Revolution (revisited):} When religion (questioned for the first time) and economy (weakened by the Seven Years' War and support for the American Revolution) both malfunctioned together, the wider social order collapsed into anarchy .
\end{itemize}
\end{CriticalPoint}

\begin{QuickRevision}
\begin{itemize}
    \item Structural functionalism = consensus theory; society = human body (organs maintain social equilibrium / homeostasis) .
    \item Pathology = organ malfunction; cure = social reform (revolution is an extreme last resort) .
    \item Coinciding malfunction of multiple organs (religion + economy) leads to social collapse/anarchy .
\end{itemize}
\end{QuickRevision}

\sectiondivider

\section{Marxism, Sociological Imagination \& Total Institutions}

\subsection{Marxism}

\begin{ComparisonBox}
\begin{itemize}
    \item \textbf{Structural Functionalists:} Hold that all parts/organs of society (religion, economy, state, family, media, law) are equally important, each contributing its own function to maintain society (e.g., economy distributes resources, investment, jobs, and growth) .
    \item \textbf{Karl Marx:} Disagrees---in the human body, one organ (the brain) controls all others and is most important . By the same logic (organic analogy), in society, the \textbf{economy is the brain} and is therefore the most important organ; all other organs are secondary and controlled by the economy .
\end{itemize}
\end{ComparisonBox}

\begin{KeyConcept}
\textbf{Marx's model of society rests on two core claims:}
\begin{enumerate}
    \item \textbf{Base and Superstructure:} Economy is the \textbf{base} of society; all other institutions (religion, law, family, state---termed the \textbf{`superstructure'}) are secondary and determined/controlled by the economic base .
    \item \textbf{Two-Class Model:} Society consists fundamentally of two classes: the \textbf{bourgeoisie} and the \textbf{proletariat} .
\end{enumerate}
\end{KeyConcept}

\paragraph{Bourgeoisie, Proletariat \& Means of Production}
\begin{itemize}
    \item Every historical epoch runs on a specific means of production . In the feudal era, land was the main source of livelihood---the dominant group controlled land (zamindars/feudal lords); the dependent group worked the land for wages (kisans/peasants) .
    \item With the Industrial Revolution, the economy shifted from agrarian (primary sector) to manufacturing (secondary sector); mass production required machines, and the dominant group became those who owned industries/factories .
    \item \textbf{Bourgeoisie:} Those who own the means of production (owners of capital/industry) .
    \item \textbf{Proletariat (working class):} Those who do not own the means of production and depend entirely on selling their labour power for wages .
\end{itemize}

\begin{ExampleBox}
\textbf{Worked example (Motorcycle design):}
\begin{itemize}
    \item A talented engineer designs a new motorcycle and sells the design to a company owner for a one-time payment of \rupee 2,00,000 .
    \item The company manufactures and sells 1,00,000 units at \rupee 80,000 each---generating a total revenue of \rupee 800 crore ($80{,}000 \times 1{,}00{,}000 = 8{,}00{,}00{,}00{,}000$) .
    \item The designer receives only \rupee 2 lakh, while the bourgeoisie owner captures the massive ongoing surplus value/profit generated by that labour .
\end{itemize}
\end{ExampleBox}

\paragraph{How the Superstructure Supports the Bourgeoisie}
\begin{itemize}
    \item \textbf{Family:} When the exploited designer complains about being underpaid relative to corporate profits, the family (spouse/parents) pacifies him (e.g., advising him to gift his boss on Diwali, or focus on job security), preventing revolt and returning him quietly to work . Family thus functions to maintain capitalist stability .
    \item \textbf{Religion:} When a worker seeks a raise or promotion and is told by the employer to read the \textit{Bhagavad Gita} (``\textit{karam karo, phal ki chinta mat karo}''---do your duty, do not worry about the fruits/rewards), religion is invoked as an ideological tool to induce patient subservience without challenging structural exploitation .
    \item \textbf{Law:} An ordinary person stealing \rupee 1,00,000 is swiftly arrested and jailed, whereas large-scale corporate fraud and corruption by economic elites rarely results in equal punishment . Economist \textbf{Kaushik Basu} (former Chief Economic Advisor, GoI) termed this the \textbf{`Sanskritisation of corruption'}---elite corruption is normalised and shielded .
    \item \textbf{State:} During the COVID-19 lockdown, multiple state governments suspended labour laws (e.g., minimum-wage guarantees, working-hour regulations under the Factory Act 1948/1961, and the Maternity Benefit Act) .
    \begin{itemize}
        \item \textit{Justification:} Attracting Foreign Direct Investment (FDI) and improving India's rank on the World Bank's Ease of Doing Business Index .
        \item \textit{Capital flight constraint:} Raising wages or labour protections prompts multinational capital to relocate to countries with cheaper labour, severely constraining the state's capacity to act purely in workers' interests .
    \end{itemize}
    \item \textbf{Conclusion (Marx):} Superstructural institutions promote values that conceal exploitation . Only when workers develop \textbf{`true consciousness'} (as opposed to false consciousness) can they unite to overthrow capitalism through revolution, establishing an egalitarian socialist/communist society where workers collectively control the means of production .
\end{itemize}

\begin{ComparisonBox}
Both structural functionalism and Marxism emerged as theoretical reactions coinciding with the rise of modernity (modern state, formal law, and nuclear family are all products of modernity) .
\end{ComparisonBox}

\begin{QuickRevision}
\begin{itemize}
    \item \textbf{Marx's model:} Economy = Base (the brain); Religion, law, family, state = Superstructure (serves the bourgeoisie) .
    \item \textbf{Classes:} Bourgeoisie (owners of means of production) vs. Proletariat (sellers of labour power) .
    \item Superstructural institutions pacify dissent, legitimate inequality, and conceal surplus-value extraction .
    \item \textbf{Transformation:} Requires moving from false consciousness to \textbf{true consciousness} $\rightarrow$ revolution $\rightarrow$ worker ownership of means of production .
\end{itemize}
\end{QuickRevision}

\subsection{Sociological Imagination --- Personal Trouble vs. Public Issue}

\begin{ExampleBox}
\begin{itemize}
    \item \textbf{Jewellery shop example:} A shop assistant earning \rupee 40,000 per month steals jewellery worth \rupee 3,00,000 . Common sense brands him a `thief' . From a Marxist perspective, the real thief is the owner who appropriates the massive surplus value produced by the assistant's labour while paying a subsistence wage .
    \item \textbf{Societal values:} A society obsessed with material success (e.g., the `American Dream' glorified through mass media and cinema) structurally propels deprived individuals toward crime; the root cause lies in social values, not individual pathology .
\end{itemize}
\end{ExampleBox}

\paragraph{Personal Trouble vs. Public Issue}

\begin{ExampleBox}
\begin{itemize}
    \item \textbf{Infertility:} One couple unable to conceive is a personal trouble (medical consultation); when millions face infertility, it becomes a public issue (declining fertility rates, state demographic policies like China's former one-child policy) .
    \item \textbf{Unemployment:} One person unemployed is a personal trouble; when 50 million people are unemployed, it is a structural public issue of economic organisation .
\end{itemize}
\end{ExampleBox}

\begin{KeyConcept}
\textbf{Sociological Imagination (C. Wright Mills):} The capacity to shift from one perspective to another---to understand that what appears to be an individual's personal trouble is rooted in public issues and structural arrangements of society .
\begin{itemize}
    \item \textbf{Critique of liberal democracy:} Formal legal equality (e.g., Article 14 of the Indian Constitution) fosters the illusion of equal opportunity, encouraging society to blame individuals for structural failures .
    \item \textbf{Suicide:} A single suicide appears as a personal tragedy; but 8--9 million suicides globally each year reflect deep structural causes . Marriage acts as an integrating institution that reduces suicide risk by embedding individuals in supportive networks .
    \item \textbf{Workplace boredom:} An individual feeling bored is a personal trouble; hundreds of millions experiencing chronic monotony reflects the public issue of \textbf{alienation} under extreme division of labour .
\end{itemize}
\end{KeyConcept}

\begin{QuickRevision}
\begin{itemize}
    \item \textbf{Sociological imagination (C. Wright Mills):} Personal troubles $\equiv$ public issues rooted in social structure .
    \item \textbf{Applications:} Unemployment, infertility / declining fertility, suicide rates, and workplace alienation .
\end{itemize}
\end{QuickRevision}

\subsection{Finding Patterns in Different Institutions --- Surveillance \& Foucault}

\begin{ComparisonBox}
Common sense views school, factory, and prison as entirely unrelated . A sociological perspective identifies a shared structural pattern: each institution contains a population of `inmates' (students, workers, prisoners) subject to hierarchy, surveillance, and disciplinary power .
\end{ComparisonBox}

\begin{KeyConcept}
\textbf{Michel Foucault and Disciplinary Power:}
\begin{itemize}
    \item \textbf{The Panopticon:} Designed by Jeremy Bentham as an architectural model for prisons (a central watchtower surrounded by peripheral cells) . Inmates can see the tower but cannot see inside it, creating the permanent illusion of being watched .
    \item \textbf{Disciplinary power vs. Sovereign power:} Historically, power was `sovereign'---spectacular, violent, and episodic (public executions, gallows) . Modern power is `disciplinary'---subtle, invisible, continuous, and internalised through surveillance (e.g., motorists slowing down at CCTV signs regardless of whether cameras are active) .
\end{itemize}
\end{KeyConcept}

\begin{ExampleBox}
In schools, the principal's office or staff room on an elevated floor with one-way vision mirrors the Panopticon watchtower, enforcing student discipline through potential visibility .
\end{ExampleBox}

\begin{QuickRevision}
\begin{itemize}
    \item \textbf{Foucault's Panopticon (Bentham):} Constant visibility/illusion of surveillance creates internalised self-discipline .
    \item \textbf{Shift in power:} Sovereign power (public spectacles of execution) $\rightarrow$ Disciplinary power (surveillance in schools, factories, prisons) .
\end{itemize}
\end{QuickRevision}

\subsection{Total Institutions (Erving Goffman)}

\begin{KeyConcept}
\textbf{Total Institutions (Erving Goffman):} Places of residence and work where a large number of like-situated individuals, cut off from the wider community for an appreciable period of time, lead an enclosed, formally administered round of life .
\begin{itemize}
    \item \textbf{Examples:} Mental asylums, boarding schools (e.g., Doon School, Mayo College, Wellington College), cult organisations (e.g., Brahma Kumaris headquarters at Mount Abu), and military cantonments/prisons .
    \item \textbf{Shared characteristics:}
    \begin{enumerate}
        \item \textit{Physical isolation:} Geographically segregated on city outskirts or remote hills .
        \item \textit{Distinct internal rules:} Autonomous rule systems governing all daily conduct .
        \item \textit{Enforced uniformity:} Stripping of personal identity markers (surrender of phones, civil clothes, watches) and imposition of standardised uniforms, schedules, and speech patterns .
    \end{enumerate}
    \item \textbf{Mortification of the self:} The systematic erosion and alteration of an individual's prior civilian identity and self-image upon entering a total institution .
\end{itemize}
\end{KeyConcept}

\begin{QuickRevision}
\begin{itemize}
    \item \textbf{Total institutions (Goffman):} Isolated, strictly regulated environments (prisons, asylums, boarding schools, cantonments) .
    \item Key outcome = \textbf{mortification of self} via enforced uniformity and isolation .
\end{itemize}
\end{QuickRevision}

\subsection{Common Sense vs. Sociology (Introduction)}

\begin{KeyConcept}
\begin{itemize}
    \item \textbf{Common sense:} Naturalised, unexamined assumptions, subjective opinions, and localized folklore .
    \item \textbf{Sociology:} A systematic body of knowledge grounded in empirical evidence, scientific methodology, data classification, and verification of validity and reliability to generate theoretical models .
    \item Distinctively, social theories often carry `-ism' labels (Marxism, Positivism, Functionalism), representing comprehensive theoretical frameworks .
\end{itemize}
\end{KeyConcept}

\paragraph{Universal Sociological Theories}
\begin{KeyConcept}
\begin{itemize}
    \item \textbf{Alienation (Marx):} Universal across sectors . A specialised PhD graduate or corporate employee becomes alienated from broader human capacities through hyper-specialisation . The IAS avoids alienation through cross-departmental administrative rotations, unlike single-domain services like the IRS . A homemaker restricted purely to domestic chores similarly suffers alienation from wider social potentials .
    \item \textbf{Reference Group Theory:} Coined by Herbert Hyman and developed by Robert K. Merton . Individuals evaluate their achievements against a chosen standard group, experience \textbf{relative deprivation}, and conform to the reference group's norms and routines (e.g., UPSC aspirants adopting the study schedules of successful toppers) .
\end{itemize}
\end{KeyConcept}

\begin{QuickRevision}
\begin{itemize}
    \item Common sense = assumption/opinion; Sociology = empirical scientific methodology .
    \item Universal theories: Marx's \textbf{Alienation} (hyper-specialisation) and Merton's \textbf{Reference Group Theory} (relative deprivation and emulation) .
\end{itemize}
\end{QuickRevision}

\sectiondivider

\section{Common Sense vs. Sociology, Scope \& Schools of Thought}

\subsection{Common Sense vs. Sociology: Studying the Same Object Differently}

\begin{ComparisonBox}
\begin{itemize}
    \item \textbf{Constable:} Executes administrative orders to apprehend an offender .
    \item \textbf{IPS Officer / Administrator:} Formulates district policies and policing strategies to reduce crime rates .
    \item \textbf{Psychologist:} Studies internal cognitive traits, deviance, and the `criminal mind' .
    \item \textbf{Sociologist:} Investigates why crime exists as a structural phenomenon, how social norms define deviance, and who possesses the power to criminalise specific acts .
    \item \textbf{Relativity of crime:} Killing another human is murder in civil society, but celebrated as valour in warfare against an enemy soldier---demonstrating that criminality is a socially constructed definition .
\end{itemize}
\end{ComparisonBox}

\paragraph{Common Sense is Local and Temporal --- Not Universal}

\begin{ExampleBox}
\begin{itemize}
    \item \textbf{Temporal variability:} Slavery was universally accepted for centuries but is now a heinous crime; women's disenfranchisement was once standard; homosexuality was criminalised historically, decriminalised in the UK/Europe around 2014, with same-sex marriage battles continuing today .
    \item \textbf{Spatial/cultural variability:} Dietary taboos against chicken on Tuesdays exist in North India but not in South or North-East India; the matrilineal \textbf{Khasi} of Meghalaya trace property through daughters; the \textbf{Mosuo} tribe of Yunnan (China) practice walking marriages with no formal marriage institution or paternal co-residence .
    \item \textbf{Sociological theory} seeks universal structural explanations across these disparate manifestations .
\end{itemize}
\end{ExampleBox}

\begin{QuickRevision}
\begin{itemize}
    \item Common sense is time-bound (temporal) and place-bound (local/particular) .
    \item Sociology analyzes how definitions of reality (e.g., crime, family, marriage) are socially constructed and historically contingent .
\end{itemize}
\end{QuickRevision}

\subsection{Comparative Method: Caste and Race}

\begin{ComparisonBox}
Common sense treats caste (ritual hierarchy) and race (biological phenotype) as completely separate categories . The comparative method reveals shared structural characteristics :
\begin{enumerate}
    \item \textbf{Endogamy:} Strict rules requiring marriage within the designated sub-group .
    \item \textbf{Commensality restrictions:} Taboos against inter-dining and physical proximity .
    \item \textbf{Segregated rules and status hierarchies:} Rigid caste/racial codes enforced by dominant groups .
\end{enumerate}
Sociologist \textbf{Oliver Cromwell Cox} famously observed that \textit{``race is nothing but coloured caste''} .
\end{ComparisonBox}

\begin{QuickRevision}
\begin{itemize}
    \item \textbf{Comparative method:} Uncovers structural parallels between caste and race (endogamy, commensal taboos, occupational closure) .
    \item \textbf{Oliver Cox:} Race is structurally equivalent to a system of coloured caste .
\end{itemize}
\end{QuickRevision}

\subsection{Functions of Institutions \& Public Issues}

\begin{KeyConcept}
Structural functionalists uncover systemic functions invisible to common-sense perception .
\begin{itemize}
    \item \textbf{Function of Law:} Functions like the biological \textbf{nervous system}---enforcing social discipline, resolving disputes, defining boundaries of deviance, and re-establishing social equilibrium .
\end{itemize}
\end{KeyConcept}

\begin{ExampleBox}
\textbf{Population Growth as a Public Issue:}
\begin{itemize}
    \item A family having five children appears as a personal choice . When aggregate data shows national expansion (e.g., India adding millions annually), it is a structural public issue .
    \item Sociological inquiry identifies the root cause not merely as `illiteracy' (a statistical correlation), but as deeply entrenched patriarchal values manifesting as \textbf{`son-meta preference'}---continuing childbearing until male heirs are born .
\end{itemize}
\end{ExampleBox}

\begin{QuickRevision}
\begin{itemize}
    \item Law acts as the nervous system of society, containing deviance .
    \item Demographic growth is structurally explained through patriarchy and son-meta preference rather than simple illiteracy .
\end{itemize}
\end{QuickRevision}

\subsection{Unpaid Labour, Data as Work \& Defamiliarisation}

\begin{CriticalPoint}
\begin{itemize}
    \item \textbf{Data as Unpaid Work:} Browsing digital platforms (Facebook, Amazon, Myntra) is experienced as leisure, but user activity generates commercial data harvested for corporate profitability---unpaid digital labour .
    \item \textbf{Unpaid Domestic Work:} Devaluing housework as non-productive is a patriarchal social construct unsupported by physical energy expenditure .
    \begin{itemize}
        \item \textbf{Selma James (1972):} Launched the \textit{`Wages for Housework'} campaign .
        \item \textbf{Policy impact:} In January 2021, Kamal Haasan's political party in Tamil Nadu included wages for housewives in its electoral manifesto, triggering widespread emulation by rival political parties .
    \end{itemize}
    \item \textbf{Defamiliarisation:} A concept formulated by \textbf{Zygmunt Bauman}---sociology acts to make the familiar strange, breaking through common-sense naturalisations to reveal underlying social construction .
\end{itemize}
\end{CriticalPoint}

\begin{QuickRevision}
\begin{itemize}
    \item Digital engagement and domestic chores represent unrecognised, unpaid labour .
    \item \textbf{Defamiliarisation (Zygmunt Bauman):} Revealing that what appears `natural' is socially constructed .
\end{itemize}
\end{QuickRevision}

\subsection{Scope of Sociology}
Sociology examines everything from interpersonal interactions to macro-historical processes .

\begin{ExampleBox}
\begin{itemize}
    \item \textbf{Micro:} Classroom theatrical role performances (Shakespeare: \textit{``All the world's a stage''}) .
    \item \textbf{Macro:} Global inequality, such as dependency theories arguing European industrial development occurred via colonial underdevelopment of India .
\end{itemize}
\end{ExampleBox}

\begin{CriticalPoint}
\begin{itemize}
    \item Sociology emerged to study the crises and transformations of modernity, leading to two core debates: (1) What is the legitimate boundary/scope of the discipline? (2) How does sociology study subjects already covered by economics or psychology? 
    \item \textit{Open question on science:} Is Yoga a science? The Indian Medical Association (IMA) often classifies AYUSH practices as pseudo-science; evaluating such claims requires first defining the rigorous criteria of scientific epistemology .
\end{itemize}
\end{CriticalPoint}

\begin{QuickRevision}
\begin{itemize}
    \item Scope ranges from micro-interaction to macro-structural dependency .
    \item Core task: Distinguishing sociological methodology from sister social sciences .
\end{itemize}
\end{QuickRevision}

\subsection{Synthetic School vs. Formalist School}

\begin{ComparisonBox}
\begin{itemize}
    \item \textbf{Synthetic School (Mainly French --- Comte, Durkheim, Hobhouse, Ginsberg):} Believed sociology must synthesize all social sciences, studying society as an integrated totality---``everything under the sun'' .
    \item \textbf{Formalist / Specialistic School (Mainly German --- Simmel, Weber, Tönnies, Von Wiese):} Argued sociology should be a distinct, specialised science studying specific forms of social interaction and human action rather than society in its entirety .
\end{itemize}
\end{ComparisonBox}

\paragraph{Durkheim's Three Branches of Sociology (Synthetic School)}
\begin{enumerate}
    \item \textbf{Social Morphology:} Demography, population density, geographical distribution, sex ratio, fertility, and mortality .
    \item \textbf{Social Physiology:} Dynamic social institutions---religion, law, economy, politics, and culture .
    \item \textbf{General Sociology:} Formulating universal sociological laws and synthesizing overarching social principles .
    \item \textit{Note:} \textbf{Social Pathology} is often recognized as a branch studying societal breakdowns, criminality, communal violence, and domestic abuse .
\end{enumerate}

\paragraph{Morris Ginsberg's Classification (Synthetic School)}
Ginsberg divided sociology into four distinct branches :
\begin{enumerate}
    \item \textbf{Social Morphology:} Population dynamics and social structures .
    \item \textbf{Social Control:} Law, religion, morals, customs, and conventions .
    \item \textbf{Social Processes:} Interactions such as cooperation, competition, conflict, and assimilation .
    \item \textbf{Social Pathology:} Maladjustments, poverty, crime, and social disintegration .
\end{enumerate}

\paragraph{Formalist Contributions: Tönnies and Simmel}
\begin{ComparisonBox}
\textbf{Ferdinand Tönnies: Typology of Human Association}
\begin{itemize}
    \item \textbf{Gemeinschaft (Community / Pre-industrial):} Characterised by primary relationships, high intimacy, emotional bonds, collective consciousness (``we-feeling''), and traditional rural life (Gandhian village model) .
    \item \textbf{Gesellschaft (Association / Post-industrial):} Characterised by contractual relations, individualism (``I-feeling''), impersonal exchange, instrumental rationality, and the urban \textbf{`blasé attitude'} (indifference toward others in urban environments) .
\end{itemize}
\end{ComparisonBox}

\begin{KeyConcept}
\textbf{Georg Simmel: Forms of Social Interaction}
Simmel argued sociology should study the \textbf{forms of association} rather than the specific content . Universal recurring forms include:
\begin{itemize}
    \item Cooperation and Conflict 
    \item Centralisation and Decentralisation 
    \item Superordination and Subordination 
\end{itemize}
\end{KeyConcept}

\begin{QuickRevision}
\begin{itemize}
    \item \textbf{Synthetic School (Comte, Durkheim):} Sociology is an encyclopedic synthesis . Durkheim's branches: Morphology, Physiology, General Sociology .
    \item \textbf{Formalist School (Simmel, Weber, Tönnies):} Specific forms of interaction .
    \item \textbf{Tönnies:} Gemeinschaft (community/we-feeling) vs. Gesellschaft (association/blasé attitude) .
    \item \textbf{Simmel:} Focus on forms (cooperation/conflict, hierarchy) rather than content .
\end{itemize}
\end{QuickRevision}

\sectiondivider

\section{Social Action Theory, Caste Dynamics \& Sociology vs. Psychology}

\subsection{Sociology Becoming Part of Common Sense}

\begin{CriticalPoint}
While sociology critiques common sense, sociological concepts diffuse into public discourse, becoming commonsensical over generations :
\begin{itemize}
    \item \textbf{`Defence mechanism'} --- Coined by Sigmund Freud, now ubiquitous in daily speech .
    \item \textbf{`Role model'} --- Formulated by Robert K. Merton, now standard vocabulary .
    \item \textbf{`Vote bank'} --- Coined by Indian sociologist M.~N. Srinivas, now central to political journalism .
    \item \textbf{Anthony Giddens} conceptualised this process as the reflexive appropriation of sociological knowledge by lay actors .
\end{itemize}
\end{CriticalPoint}

\begin{QuickRevision}
\begin{itemize}
    \item Concepts coined in sociology/psychology (defence mechanism, role model, vote bank) regularly pass into everyday common sense (Anthony Giddens) .
\end{itemize}
\end{QuickRevision}

\subsection{Caste: Consensus and Conflict Perspectives}

\begin{GovtPolicy}
\textbf{Decentralisation and Conflict (Simmel's Forms):}
The \textbf{Hindu Succession (Amendment) Act, 2005} granted equal coparcenary property rights to daughters . This legal decentralisation of inheritance from sons-only to sons-and-daughters disrupted traditional patriarchal centralisation, creating new intra-familial conflicts over property despite customary familial affection .
\end{GovtPolicy}

\begin{ComparisonBox}
\textbf{Dual Nature of Caste Organisation:}
\begin{itemize}
    \item \textbf{Consensus / Functional Model:} The traditional \textit{Chaturvarna} division of labour (Brahmins = knowledge, Kshatriyas = protection, Vaishyas = trade, Shudras = service) functioned cooperatively without direct economic competition . Early Gandhi highlighted this functional interdependence .
    \item \textbf{Conflict Model:} Rigid caste stratification enforces severe systemic oppression, untouchability, commensal bans, and temple-entry denials (historically against Dalits; gendered exclusions as in the Sabarimala dispute) . Gandhi's perspective shifted toward anti-untouchability following B.~R. Ambedkar's assertions around the Poona Pact (1932) .
\end{itemize}
\end{ComparisonBox}

\begin{QuickRevision}
\begin{itemize}
    \item Hindu Succession Amendment illustrates how institutional decentralisation generates new relational conflict .
    \item Caste can be viewed through consensus (non-competitive division of labour) or conflict (structural dominance and ritual exclusion) .
\end{itemize}
\end{QuickRevision}

\subsection{Institutions in Transition: Gemeinschaft to Gesellschaft}

\begin{ComparisonBox}
\begin{itemize}
    \item \textbf{Popular Culture:} 1990s Bollywood cinema portrayed multi-generational joint families led by patriarchal elders (Gemeinschaft); contemporary cinema features individualistic, career-driven protagonists (Gesellschaft) .
    \item \textbf{Gender Relations:} Subjugated notions of female obligation (``\textit{Dharam Patni}'') have given way to individual rights, legal autonomy, and collective resistance (the \textbf{\#MeToo} movement representing the shift from `we' to `me/I') .
    \item \textbf{Family Formations:} Emergence of \textbf{`Living Apart Together' (LAT)}---couples maintaining separate domestic residences due to dual professional careers (e.g., civil servant couples in different postings) .
\end{itemize}
\end{ComparisonBox}

\begin{QuickRevision}
\begin{itemize}
    \item The transition to Gesellschaft manifests in delayed marriages, rising divorce rates, individual careerism, and non-traditional family structures .
\end{itemize}
\end{QuickRevision}

\subsection{Max Weber and Social Action Theory}

\begin{KeyConcept}
\textbf{Social Action (Max Weber --- Formalist School):} Action is `social' insofar as its subjective meaning takes account of the behaviour of others and is thereby oriented in its course .
\begin{itemize}
    \item Structural functionalism and Marxism analyze macro-structures; Weber focuses on the micro-level \textbf{subjective meanings and motives} of individual actors .
    \item \textit{Institutional contrast:} Functionalism sees marriage as ensuring societal reproduction; Weber examines individual motives (e.g., securing permanent residency/visas, emotional intimacy, or financial security) .
\end{itemize}
\end{KeyConcept}

\begin{CriticalPoint}
\textbf{The Hathras Case \& Religious Conversion as Social Action:}
\begin{itemize}
    \item Following the Hathras atrocity (the gang rape of a Dalit woman by upper-caste men), affected Dalit villagers converted to Buddhism as a deliberate act of political and religious protest .
    \item Functionalism cannot explain this phenomenon (since it views religion purely as a source of social harmony/order); conflict theory and social action theory explain it as an assertion of dignity and resistance against caste hierarchy .
    \item \textbf{Intersectionality:} Dalit women face compound oppression across the dual axes of caste hierarchy and gender subjugation .
    \item \textbf{Dr. B.~R. Ambedkar:} Challenged the theological foundations of Hinduism (author of \textit{Buddha and Karl Marx}); established the \textbf{Navayana} (`new vehicle') school of Buddhism, rejecting traditional Hinayana/Mahayana/Vajrayana quietism and redefining spiritual liberation as social emancipation .
\end{itemize}
\end{CriticalPoint}

\paragraph{Macro, Meso, and Micro Levels of Sociological Analysis}

\begin{KeyConcept}
\begin{enumerate}
    \item \textbf{Macro-Sociology:} Grand structures, total systems, economy, state, and global patterns (Structural Functionalism, Marxism) .
    \item \textbf{Meso-Sociology:} Intermediate institutions, formal organisations, communities, and forms of relationship (Simmel, Tönnies) .
    \item \textbf{Micro-Sociology:} Individual social action, face-to-face interaction, and subjective meanings (Weber, Goffman, Mead) .
\end{enumerate}
\end{KeyConcept}

\begin{QuickRevision}
\begin{itemize}
    \item \textbf{Max Weber:} Social action focuses on micro-level subjective meanings and motives .
    \item \textbf{Ambedkar \& Navayana:} Reinterpretation of Buddhism as emancipatory social action against caste .
    \item Sociological scope operates across \textbf{Macro} (structures), \textbf{Meso} (groups/forms), and \textbf{Micro} (action/meanings) .
\end{itemize}
\end{QuickRevision}

\subsection{Sociology vs. Psychology}

\begin{ComparisonBox}
\begin{itemize}
    \item \textbf{Psychology:} Studies individual cognition, mental states, neuro-biological systems, emotions, personality development, and internal pathologies in isolation from wider social structures .
    \item \textbf{Sociology:} Analyzes how social structures, institutions, and collective norms shape individual action and consciousness (examines the public/structural roots of personal problems) .
\end{itemize}
\end{ComparisonBox}

\paragraph{Freudian Psychoanalysis vs. Sociological Self}
\begin{KeyConcept}
\begin{itemize}
    \item \textbf{Sigmund Freud:} The psyche comprises the \textbf{Id} (instinctual, irrational, aggressive impulses), \textbf{Ego} (rational executive mediator), and \textbf{Superego} (internalised moral conscience) . Investigated via psychoanalysis (free association, dream analysis) . Freud viewed personality development as an internal psychic struggle largely independent of wider social variation .
    \item \textbf{Sociological Response (The Social Self):} George Herbert Mead demonstrated that the self develops socially through language and role-taking:
    $$\text{Self} = \textbf{`I'} \text{ (spontaneous, impulsive ego)} + \textbf{`Me'} \text{ (socialised, internalised expectations of society)}$$
    Consumer choices (e.g., buying an iPhone for social validation) reflect the socially conditioned `Me' .
\end{itemize}
\end{KeyConcept}

\paragraph{Theoretical Bridges: Roles and Social Psychology}
\begin{KeyConcept}
\begin{itemize}
    \item \textbf{Roles as building blocks:} Individuals do not act in a vacuum; they enact socially defined \textbf{roles} (doctor, parent, student, civil servant) .
    \item \textbf{Role Conflict:} Incompatible demands arising from holding two or more distinct roles simultaneously (e.g., an army officer balancing frontline military deployment with parental care; working a full-time job while preparing for UPSC) .
    \item \textbf{Role Strain:} Stress experienced within a single role due to conflicting or excessive expectations (e.g., Sachin Tendulkar facing overwhelming pressure while pursuing his 100th international century) .
    \item \textbf{Social Psychology:} The interdisciplinary bridge investigating how individual mental functioning is conditioned by group dynamics .
\end{itemize}
\end{KeyConcept}

\begin{QuickRevision}
\begin{itemize}
    \item \textbf{Psychology} focuses on internal mental dynamics; \textbf{Sociology} focuses on social structures and context .
    \item Freud's Id/Ego/Superego countered by Mead's Social Self (`I' and `Me') .
    \item Bridges: Social action theory, \textbf{role conflict} (between roles), and \textbf{role strain} (within a single role) .
\end{itemize}
\end{QuickRevision}

\sectiondivider

\section{Suicide, Impression Management, History \& Economics}

\subsection{Suicide: Psychological vs. Sociological Perspective (Durkheim)}

\begin{ComparisonBox}
\begin{itemize}
    \item \textbf{Psychological approach:} Investigates idiosyncratic clinical factors---mental depression, neurosis, chemical imbalances, or private distress .
    \item \textbf{Sociological approach (\'Emile Durkheim):} Disregards individual psychological profiles to examine aggregate suicide rates across distinct social groups and institutions .
\end{itemize}
\end{ComparisonBox}

\paragraph{Durkheimian Demographic Patterns}
Durkheim's empirical study revealed systematic group-level regularities :
\begin{itemize}
    \item \textbf{Unmarried vs. Married:} Single individuals exhibit higher suicide rates than married individuals .
    \item \textbf{Gender:} Men display higher suicide rates than women .
    \item \textbf{Peace vs. War:} Suicide rates consistently decline during periods of national wartime (e.g., during the Kargil War) and rise during peacetime .
    \item \textbf{Religion:} Protestants display higher suicide rates than Catholics due to differences in ecclesiastical integration .
\end{itemize}

\begin{KeyConcept}
\textbf{Social Integration as an Explanatory Variable:}
Suicide rates vary inversely with the degree of social integration . Marriage binds the individual into a reciprocal web of domestic obligations; wartime generates intense collective solidarity and patriotism against a common adversary, submerging individual despair into national cohesion .
\end{KeyConcept}

\begin{QuickRevision}
\begin{itemize}
    \item Durkheim proved suicide is a social fact governed by degrees of \textbf{social integration} and \textbf{social regulation} .
    \item Higher rates among single individuals, men, and peacetime societies .
\end{itemize}
\end{QuickRevision}

\subsection{`Individuals Do Not Exist, Only Roles Exist' \& Impression Management}

\begin{KeyConcept}
\textbf{Dramaturgical Perspective (Erving Goffman):}
Social life is staged like a theatrical performance where individuals manage impressions to project desired identities .
\begin{itemize}
    \item \textbf{Front Stage:} Formal social settings where individuals perform before an audience according to strict scripts and role expectations (e.g., lecturing in class, meeting clients, answering a phone call from parents while outside) .
    \item \textbf{Backstage:} Private spaces where performers can drop their public facade, relax, express hidden emotions, and rehearse future performances (e.g., lounging alone at home) .
    \item \textbf{Impression Management:} The deliberate cognitive and behavioural adjustments made by an individual to control the information received by others in social interactions .
\end{itemize}
\end{KeyConcept}

\begin{QuickRevision}
\begin{itemize}
    \item \textbf{Goffman's Dramaturgy:} Social life = performance; Front stage (public performance) vs. Backstage (private relaxation) .
    \item \textbf{Impression management:} Continual tailoring of behaviour across social situations .
\end{itemize}
\end{QuickRevision}

\subsection{Sociology vs. History}

\begin{ComparisonBox}
\begin{itemize}
    \item \textbf{Traditional History:} Chronological, descriptive documentation of unique past events, political regimes, treaties, and individual leaders (ancient, medieval, modern) .
    \item \textbf{Sociology:} Analytical, generalizing study of modern social structures, institutional evolution, patterns of human association, and recurring structural transformations .
\end{itemize}
\end{ComparisonBox}

\paragraph{Competing Historiographies of the French Revolution}
Historical facts do not speak for themselves; theoretical perspectives determine their meaning :
\begin{itemize}
    \item \textbf{Marxist Historiography (Georges Lefebvre, Karl Marx):} The revolution was a classic bourgeois transformation where the rising capitalist class overturned feudal nobility, using popular peasant discontent to seize state power .
    \item \textbf{Revisionist Historiography (Alfred Cobban):} Challenges the Marxist thesis, demonstrating that the revolutionary leadership comprised educated professionals, jurists, and civil servants rather than industrial capitalists .
    \item \textbf{Conspiracy / Intellectualist View:} Frames the revolution as an orchestrated campaign by Enlightenment intellectuals against Church hegemony .
\end{itemize}

\paragraph{Feminist Historiography}
\begin{CriticalPoint}
Mainstream nationalist historical narratives focus almost exclusively on male figures (Gandhi, Nehru, Patel, Ambedkar) while obscuring women's mass mobilization .
\begin{itemize}
    \item \textbf{Dandi Salt March (1930):} Mahatma Gandhi initially opposed women's active participation in the march; women were included only following the fierce intervention of \textbf{Sarojini Naidu} . Historical interpretations diverge on whether Gandhi sought to protect women from police brutality or reflected patriarchal assumptions about female capability .
\end{itemize}
\end{CriticalPoint}

\paragraph{Vertical vs. Horizontal Mobilisation}
\begin{ComparisonBox}
\begin{itemize}
    \item \textbf{Vertical Mobilisation:} An elite, Western-educated upper-caste leader mobilising subaltern masses from above (e.g., Gandhi during the Champaran Satyagraha of 1917)---which can inadvertently reinforce the assumption that subalterns require elite patronage .
    \item \textbf{Horizontal Mobilisation:} Subaltern groups mobilising autonomously under their own leadership to assert structural rights (e.g., \textbf{Periyar E.~V. Ramasamy's Self-Respect Movement} in Tamil Nadu against Brahminical ritual dominance) .
\end{itemize}
\textbf{Structural Representation Data:} Upper castes continue to hold disproportionate representation in higher bureaucracy and Parliament ($\sim 80\%$ of nominated Rajya Sabha members historically) . Scholar \textbf{Kancha Ilaiah} polemically termed the IAS the \textit{`Indian Brahminical Service'} .
\end{ComparisonBox}

\paragraph{Subaltern, Colonialist, Orientalist, and Nationalist Historiography}
\begin{KeyConcept}
\begin{itemize}
    \item \textbf{Subaltern Historiography (Ranajit Guha, Shahid Amin):} Reconstructs history from the perspective of marginalized social groups (peasants, tribals, women, Dalits) ignored by elite historiography .
    \begin{itemize}
        \item \textbf{Shahid Amin (Gorakhpur, 1921--22):} Analyzed the `making of the Mahatma', showing that local peasant myths, religious visions, and rumours transformed Gandhi into a messianic figure, challenging the simplistic narrative that the title was merely conferred by Rabindranath Tagore .
        \item \textbf{Literary Subaltern Voices:} Mulk Raj Anand (\textit{Coolie}, \textit{Untouchable}) and Arundhati Roy (\textit{The God of Small Things}) center subaltern experiences; R.~K. Narayan portrays everyday life .
    \end{itemize}
    \item \textbf{Colonialist Historiography (James Mill):} Periodised Indian history into Hindu, Muslim, and British eras, presenting colonial rule as a `civilising mission' .
    \item \textbf{Orientalist Historiography (William Jones, Max Müller):} Celebrated ancient Sanskrit texts and Indian classical philosophy (founding of the Asiatic Society) .
    \item \textbf{Nationalist Historiography (Jawaharlal Nehru --- \textit{Discovery of India}):} Emphasised civilisational unity, composite culture, and anti-colonial resistance .
\end{itemize}
\end{KeyConcept}

\begin{KeyConcept}
\textbf{Mutual Dependence of Sociology and History:}
\textit{``Sociology is rootless without history; history is fruitless without sociology''} (formulated by Charles Tilly) .
\begin{itemize}
    \item \textbf{Historical Sociology:} Examines institutional trajectories over time (e.g., how caste transitioned from an occupational hierarchy into a vehicle for modern democratic political mobilization) .
    \item \textbf{Comparative Revolution Theory (Theda Skocpol):} Whereas historians focus on the differences between the French, Russian, Chinese, and American revolutions, Skocpol's \textit{States and Social Revolutions} identified a common structural causality: state breakdown caused by the coincidence of \textbf{international geopolitical/financial pressure} and \textbf{domestic peasant revolt} .
\end{itemize}
\end{KeyConcept}

\begin{QuickRevision}
\begin{itemize}
    \item Historiography proves facts require theoretical interpretation (Marxist, Revisionist, Feminist, Subaltern, Colonialist, Nationalist) .
    \item \textbf{Tilly:} Sociology is rootless without history; history is fruitless without sociology .
    \item \textbf{Theda Skocpol:} Common structural patterns of revolution (external pressure + internal revolt) .
\end{itemize}
\end{QuickRevision}

\subsection{Sociology vs. Economics}

\begin{ComparisonBox}
\begin{itemize}
    \item \textbf{Classical Economics:} Assumes markets are autonomous, self-regulating systems operating via pure price mechanisms of supply and demand, `disembedded' from social relations .
    \item \textbf{Sociology:} Assumes economic actions are deeply \textbf{embedded} in social networks, political structures, cultural values, and state institutions .
    \item \textbf{Empirical instances of state market regulation:} The \textit{Pradhan Mantri Garib Kalyan Anna Yojana} during the pandemic, post-2008 regulatory interventions across global capital markets, and the Reserve Bank of India's regulatory frameworks managing non-performing assets (NPAs) .
\end{itemize}
\end{ComparisonBox}

\paragraph{Homo Economicus vs. Bounded Rationality}
\begin{KeyConcept}
\begin{itemize}
    \item \textbf{Homo Economicus:} Classical economic model of man as a perfectly rational, atomistic, profit-maximizing, utility-calculating actor .
    \item \textbf{Sociological Limitations:} Cannot explain altruism, charity, voluntary sacrifice, unpaid care work, or caste-based occupational stratification where lower castes remain clustered in hazardous occupations despite formal legal market access .
    \item \textbf{Veblen Goods \& Conspicuous Consumption (Thorstein Veblen):} Goods purchased primarily to signal wealth and social prestige (e.g., luxury iPhones, high-end watches, five-star hospitality), directly violating the economic law of demand .
    \item \textbf{Bounded Rationality:} Economic decisions are framed, bounded, and directed by social norms, peer pressure, and institutional contexts .
\end{itemize}
\end{KeyConcept}

\paragraph{Max Weber: The Protestant Ethic and the Spirit of Capitalism}
\begin{KeyConcept}
Where economics explains capitalism through capital accumulation and technical innovations, Max Weber identified religious ideas as a primary catalyst :
\begin{itemize}
    \item Weber demonstrated that modern rational capitalism emerged predominantly in regions dominated by \textbf{Calvinist / Protestant asceticism} .
    \item \textbf{Maxims of Benjamin Franklin:} \textit{``Time is money''} and \textit{``Work is worship''} encapsulated a religiously sanctioned work ethic .
    \item Diligent labour in one's worldly calling (the `calling') and systematic reinvestment of wealth---rather than lavish indulgence---functioned as psychological signs of divine salvation (doctrine of predestination), laying the cultural foundations for modern industrial capitalism .
\end{itemize}
\end{KeyConcept}

\begin{QuickRevision}
\begin{itemize}
    \item \textbf{Embeddedness:} Economy is situated within social, cultural, and political institutions (critique of disembedded self-regulating markets) .
    \item \textbf{Veblen Goods \& Conspicuous Consumption:} Status-signaling violates classical price-demand mechanics .
    \item \textbf{Weber's Protestant Ethic:} Protestant asceticism provided the ethical ethos that stimulated modern industrial capitalism .
\end{itemize}
\end{QuickRevision}